%
%
%
%
%
%
% Alternative Chapter: "Discussion"
\chapter{Discussion: Mechanizing Semantic Proofs — Findings and Insights}

This chapter evaluates the mechanization of semantic proofs in Lean 4 by synthesizing the case-specific insights from the method chapter. The thesis examines how pen-and-paper proofs must be transformed when formalized, with particular attention to implicit assumptions that become explicit, additional machinery required, and lessons learned throughout the mechanization process.

%
%
%
\section{From Informal to Formal: The Transformation of Semantic Definitions}
\label{discussion:sec:transformation}

A primary theme emerging from the formalization of syntactic and semantic
definitions is the systematic conversion of implicit conventions to explicit
requirements.

\subsection*{Syntactic Categories: From BNF to Inductive Types}

The translation of BNF notation to Lean's inductive types is largely
mechanical but reveals several implicit assumptions in textbook presentations:

\begin{itemize}
    \item \textbf{Constructive Numerals}:
    The textbook defines numerals as binary numbers via inductive definition;
    however, this mechanization choice was not strictly necessary, as we could have
    used Lean 4's built-in \texttt{Nat} type instead. We chose the inductive
    definition to match the textbook presentation and to enable structural induction
    over the binary representation.

    \item \textbf{Naming and Scoping}:
    Informal notation permits clash-free
    use of keywords like \texttt{if} and \texttt{while}. Mechanization
    requires renaming (\texttt{cond}, \texttt{loop}) to avoid collisions
    with the host language.

    \item \textbf{Boolean Realism}:
    Choosing between symbolic notation and
    computational booleans (using Lean's native \texttt{true}/\texttt{false})
    affects how theorems about boolean logic can be leveraged, reducing
    proof effort.
\end{itemize}

\subsection*{Semantic Functions: Totality and Exhaustiveness}

Formalizing evaluation functions \(\mathcal{N}\), \(\mathcal{A}\), and
\(\mathcal{B}\) reveals a key design choice: Lean enforces total function definitions
via exhaustive pattern matching\missingcitation \todo{Cite the Lean 4 manual or a paper on Dependent Type Theory regarding totality requirements in functional languages.}, whereas textbooks define these as inductive relations
and prove totality separately. We chose to define them directly as total functions,
avoiding redundant complexity and enabling direct computation. This requires explicitly
handling all constructor cases—including those ruled out by side conditions, which
must be shown to lead to contradictions.

%
%
%
\section{The Explicit Nature of Inductive Proofs}
\label{discussion:sec:inductive_proofs}

In textbook arguments, appeals to structure (``the shape of the derivation tree'')
must be formalized as explicit case analysis on inductive constructors via Lean's
\texttt{induction} tactic, with automatically generated cases corresponding to rule labels.

\subsection*{Exhaustive Case Analysis and Contradiction Handling}

When conditions narrow which rules apply, mechanization requires explicitly deriving
contradictions. For instance, if \(\mathcal{B}[\![b]\!]s = \mathbf{tt}\) holds,
attempting to apply \texttt{if\_false} (which requires the condition false) immediately
generates a contradiction (e.g., \texttt{rw [h\_cond\_true] at h\_cond\_false; contradiction}).
This explicit handling of ruled-out cases is invisible in textbook proofs but mandatory in mechanization.

\subsection*{Modularity Through Lemmatization}

The while-loop unrolling proof demonstrates that splitting a biconditional
proof into two directed implications (\texttt{while\_expand} and
\texttt{if\_compact}) mirrors good mathematical practice. In mechanization:

\begin{itemize}
    \item \textbf{Independent Lemmas}:
    Proving $A \implies B$ and $B \implies A$ separately, then combining via
    \texttt{Iff.intro}, mirrors the textbook's two-part structure.
    \item \textbf{Constructive Proof Objects}:
    Building derivation trees explicitly (via \texttt{have} statements) enables reorganizing
    the proof steps to match the textbook's narrative structure. Rather than purely a clarity tool,
    \texttt{have} allows proof steps to be introduced earlier than their logical dependencies
    would naturally suggest, facilitating a presentation that mirrors the textbook's order.
    This can also aid reader comprehension by making the proof intent explicit.
\end{itemize}

%
%
%
\section{Multi-Step Reasoning and the Abstraction Layers}
\label{discussion:sec:multistep}

A subtle but pervasive difference emerges when formalizing proofs involving
multi-step derivations in structural operational semantics.

\subsection*{From Implicit to Explicit Multi-Step Relations}

The textbook introduces shorthand notation for multiple steps
($\Rightarrow^k$ for $k$ steps, $\Rightarrow^*$ for arbitrary steps).
Mechanization requires:

\begin{itemize}
    \item \textbf{Inductive Types for Relations}:
    The notation $\Rightarrow^k$ must be formalized as an inductive type
    \texttt{small\_step\_k} that counts steps explicitly.
    
    \item \textbf{Strong Induction on Length}:
    Proofs by ``induction on the length of the derivation sequence'' become
    \texttt{Nat.strongRecOn} in Lean, enabling strong induction where the hypothesis
    applies to any smaller step count.
    
    \item \textbf{Composition of Derivations}:
    Informal chaining becomes explicit calls to transitivity lemmas like
    \texttt{small\_step\_star\_trans}.
\end{itemize}

\subsection*{Case Study: Composition Decomposition (Lemma 2.19)}

The proof that a composition can be split into two independent executions
illustrates how mechanization requires explicit generalization, arithmetic
justification via \texttt{linarith}, and careful unpacking of existential witnesses
from induction hypotheses.

%
%
%
\section{Equivalence Proofs: Bridging Two Semantic Styles}
\label{discussion:sec:equivalence}

The two-lemma proof of equivalence between natural semantics and
structural operational semantics demonstrates how mechanization handles the relationship
between partial and total representations.

\subsection*{Lemma 2.27: Natural Semantics to Structural Semantics}

This proof constructs multi-step SOS derivations from big-step NS derivations via
induction on the shape of the NS derivation tree (pattern matching on constructors like
\texttt{ass}, \texttt{skip}, \texttt{comp}, \texttt{if\_true}, \texttt{while\_true}, etc.).
Textbook appeals to ``Exercise 2.21'' become explicit function calls; building multi-step
derivations explicitly requires knowing when to use \texttt{small\_step\_star.step} versus
\texttt{small\_step\_star\_trans}.

\subsection*{Lemma 2.28: Structural Semantics to Natural Semantics}

This proof uses strong induction on step count, requiring careful handling of
composition rules (splitting [comp$_{\text{sos}}^1$] and [comp$_{\text{sos}}^2$] into separate cases)
and explicit decomposition via Lemma 2.19. Arithmetic facts require \texttt{linarith};
the while loop unrolls into a conditional-composition structure via explicit calls to \texttt{lemma\_2\_5}.

\subsection*{Relational vs. Functional Interpretation}

While the textbook frames equivalence as ``the semantic functions are equal'',
mechanization works with evaluation relations. Lean can represent partial functions via dependent types\missingcitation \todo{Cite a source on modeling partiality in Coq/Lean/Agda, such as 'Software Foundations' or a paper on the Option/Part monad.},
but reasoning about them introduces an extra layer of proof obligation. Since equivalence proofs
must establish relational correspondence anyway, working directly with relations as primitives
is both simpler and sufficient.

%
%
%
\section{Synthesis: Core Mechanization Transformations}
\label{discussion:sec:synthesis}

The preceding sections demonstrate that mechanization is not merely a formal restatement
of textbook arguments, but a systematic transformation. The core differences fall into
several recurring patterns:

\begin{itemize}
    \item \textbf{Totality Enforcement}: Functions must be total
    (\autoref{discussion:sec:transformation}), requiring exhaustive pattern matching
    that textbooks elide.
    
    \item \textbf{Explicit Case Analysis}: Appeals to ``the shape of the derivation
    tree'' become explicit \texttt{induction} and \texttt{cases} tactics
    (\autoref{discussion:sec:inductive_proofs}), automatically generating cases for
    each constructor.
    
    \item \textbf{Contradiction Handling}: Cases ruled out by side conditions must be
    explicitly shown to lead to contradictions
    (\autoref{discussion:sec:inductive_proofs}).
    
    \item \textbf{Multi-Step Composition}: Informal chaining in SOS proofs becomes
    explicit transitivity lemmas (\autoref{discussion:sec:multistep}).
    
    \item \textbf{Generalization Strategy}: Variables that should not be fixed during
    induction must be explicitly generalized (\autoref{discussion:sec:multistep}),
    a concern invisible in informal mathematics\missingcitation \todo{Optional: Cite a paper on the 'De Bruijn factor' or general formalization overhead to support the claim that generalization is a specific 'formal' burden.}.
    
    \item \textbf{Relational vs. Functional Reasoning}: While partial functions can
    be modeled in Lean via dependent types, working directly with evaluation relations
    as primitives is both simpler and sufficient (\autoref{discussion:sec:equivalence}).
    
    \item \textbf{Automation Over Arithmetic}: Facts dismissed as ``obvious'' in the
    textbook require explicit \texttt{linarith}, \texttt{simp}, or manual case
    analysis (\autoref{discussion:sec:multistep}).
\end{itemize}

Across all these transformations, a unifying theme emerges: pen-and-paper proofs rely on
informal conventions and reader understanding that mechanization must make explicit and verifiable.


%
%
%
\section{Lessons for Future Mechanization}
\label{discussion:sec:lessons}

Based on the mechanization experience, several lessons emerge for future work on semantic formalization.

\subsection*{Planning and Structuring Proofs}

\begin{itemize}
    \item \textbf{Identify Proof Patterns Early}:
    Recognize whether the proof is structural (shape of derivations via \texttt{induction})
    or arithmetic (length of sequences via \texttt{Nat.strongRecOn}).
    
    \item \textbf{Lemmatize Aggressively}:
    Break large proofs into smaller, reusable lemmas. What appears as
    informal appeal in textbooks should become explicit lemma calls.
    
    \item \textbf{Generalization and Witnesses}:
    When setting up induction, explicitly generalize variables that should not be
    fixed. Plan how to unpack and specialize existential witnesses via pattern matching.
\end{itemize}

\subsection*{Handling Side Conditions and Automation}

\begin{itemize}
    \item \textbf{Enumerate Side Conditions}:
    Cases ruled out by side conditions must be shown to lead to contradictions
    via explicit \texttt{rw} and \texttt{contradiction} tactics.
    
    \item \textbf{Rely on Automation Where Possible}:
    Use \texttt{simp} and \texttt{linarith} liberally for routine bookkeeping,
    but separate non-arithmetic reasoning from arithmetic reasoning for clarity.
\end{itemize}

\subsection*{Relations Over Partial Functions}

Work with evaluation relations as the primary object rather than partial functions.
Partial functions cannot be directly formalized in dependent type theory\missingcitation \todo{Cite the 'halting problem' context or specific constraints of the Calculus of Inductive Constructions regarding non-termination/partiality.}; relations
are more tractable and sufficient for equivalence proofs.